%=====================================================================
% response/patches/mouse_satellite_caveat.tex
%
% WHAT THIS REPLACES
%   Three places assert a pachytene peak. Verbatim excerpts, all located by
%   `grep -n "pachytene" submission/paper.tex`:
%
%   (R1) submission/paper.tex:306--312, Results, first sentence of the mouse
%        paragraph after "(Fig.~\ref{fig:fig4}C).":
%   | The correlation rises as chromosomes
%   | synapse and peaks at pachytene, when the synaptonemal complex is fully assembled, then falls
%   | at diplotene as it disassembles: sequence similarity translates into ectopic proximity
%   | precisely during the window in which crossovers form. The stage profile is
%   | resolution-dependent---at 50 kb windows it is flat ($\rho = 0.29$--0.34)---so we report the
%   | 20 kb profile and treat the stage ordering as descriptive rather than as a tested difference.
%
%   (R2) submission/paper.tex:81--83, summary paragraph:
%   | Mouse contact maps show the same preference:
%   | the association strengthens as chromosomes synapse and peaks at pachytene, when
%   | the synaptonemal complex is fully assembled.
%
%   (R3) submission/paper.tex:545--547, Fig. 4c caption tail:
%   | The arm-level Mantel test on
%   | the same data, which respects that non-independence, gives $\rho = 0.41$, 0.58, 0.63 and
%   | 0.33 across the four stages (Methods), peaking at pachytene.}
%
% WHAT IT INSERTS
%   A satellite-aware replacement for each of the three, in decreasing length.
%   The Results replacement is the strongest wording I can defend; below it is
%   a shorter fallback for the case where the stage ordering turns out to be
%   indistinguishable after masking (the paper currently treats it as
%   descriptive anyway, so the fallback is the smaller edit).
%
% EVIDENCE PATHS, and what I re-derived from each rather than copying.
%   E1  response/notes/M6_mouse_satellite.md
%         :24--36  annotation vs proxy table; mean minor-proxy fraction 0.00004
%                  (p-end) and 0.00112 (q-end)
%         :53--66  sat_frac_union p-end min 0.2597 max 0.9928, 38/39 >= 0.50;
%                  q-end min 0.0021 max 0.0077
%         :90--96  baseline rho 0.4074/0.5785/0.6259/0.3282, exact reproduction
%                  of scripts/hic/mantel_mouse_d_m5.json
%         :112--130 3a low-satellite (21 q-q pairs) 20 kb profile
%         :160--172 "the stage structure does not survive"; peak-stage flips
%         :350--369 one-paragraph verdict
%   E2  scripts/hic/mouse_satellite_sensitivity.json -- I recomputed these myself
%       from `satellite_annotation.per_flank` (78 flanks = 39 p-end + 39 q-end),
%       not from the note:
%         p-end sat_frac_union: min 0.2597, max 0.9928, mean 0.8848,
%                               38 of 39 >= 0.50, mean frac_bins_sat50 0.8607
%         q-end sat_frac_union: min 0.0021, max 0.0077, mean 0.0044
%         p-end sat_frac_major_GSAT_MM: mean 0.2886 (min 0.0000, max 0.9720)
%         p-end sat_frac_minor_PROXY:  mean 0.00004 ; q-end mean 0.00112
%         `satellite_annotation.families` = GSAT_MM 471, MurSAT1 120,
%                                MMSAT4 168, CENSAT_MC 59, SYNREP_MM 266,
%                                IMPB_01 2548, ZP3AR 212, CYRA11_Mm 617
%         `verdicts.peak_stage_20kb_baseline` = "pachytene";
%         `verdicts.peak_stage_20kb_low_sat(both<0.25)` = "zygotene";
%         `verdicts.peak_stage_20kb_high_sat(both>=0.25)` = "zygotene"
%   E3  end-to-end-report/report/05_hic_validation.md:176 (mouse section pointer)
%         and end-to-end-report/report/08_mouse.md for the stage-Mantel framing.
%
% THREE THINGS I DELIBERATELY DO NOT CLAIM, because the evidence does not carry
% them, and each is a place the obvious sentence would be false:
%
%   1. "the p-arm satellite fraction is major satellite". It is not, alone.
%      GSAT_MM accounts for a mean 0.2886 of the p-end flank; the union mean is
%      0.8848 and the union is eight Dfam families. The minor classes are a
%      PROXY (there is no Mus "minor satellite" family in this Dfam 3.8 library:
%      MurSAT1+MMSAT4+CENSAT_MC were substituted) and recover a mean of 4e-5,
%      so minor-satellite content of these windows is UNMEASURED, not small.
%      The text therefore says what was measured and names the proxy.
%
%   2. "masking the satellite removes the association" -- no masking was done.
%      §5a of E1 and `not_established` in E2 both record that the contact axis
%      cannot be satellite-masked (needs zuo2021_<stage>.mcool plus the
%      haplotype map) and that the sequence axis cannot be re-derived without
%      rebuilding the pggb graph. What exists is a STRATIFICATION of arm pairs
%      by flank satellite fraction. The text says "setting aside the
%      satellite-rich windows", never "masked".
%
%   3. "the satellite-poor subset confirms the ordering". It does not: the
%      20 kb profile there is flat, and its nominal top stage is zygotene by
%      0.0026, a margin smaller than the Monte-Carlo noise of the 200-permutation
%      null (E1 §5c). Furthermore the so-called satellite-poor subset IS the
%      q-q subset -- 21 pairs among 7 q arms -- because arm-level satellite
%      fraction is bimodal with nothing between 0.10 and 0.25 (E1 §3). It is
%      not an independent sample and the text says so.
%
%   Also note the same flip happens in the satellite-RICH subset (zygotene
%   0.6515 vs pachytene 0.6394, E1 §3a), so "pachytene" was never robust to
%   stratification at all. This is why the replacement does not simply swap
%   pachytene for zygotene: swapping one unstable rank for another would be the
%   same error in the opposite direction.
%=====================================================================

%%% BEGIN INSERT (R1) -- Results, replaces submission/paper.tex:306--312
Because mouse chromosomes are telocentric, the 1\,Mb p-arm flanks that carry the
mouse PHRs are 26--99\% satellite (mean 88.5\% over 39 p-end flanks, 38 of them
above 50\%; all 39 q-end flanks are below 0.8\%), and satellite arrays are
near-identical between chromosomes, so these windows sit at the top of the
similarity axis and at the extreme of the multi-mapping problem at the same
time. Setting the satellite-rich windows aside, the sequence--contact
association persists over the 21 satellite-poor arm pairs ($\rho = 0.486$--$0.539$
at 20\,kb, permutation $p \leq 0.025$) and is therefore not created by
satellite-rich p-arm flanks; the stage ordering, however, does not survive the
same restriction -- the 20\,kb profile is flat across the four stages
($\rho = 0.536$, 0.539, 0.486 and 0.486) with no resolvable maximum. The
baseline 20\,kb maximum is itself fragile: the same flip of the nominal top
stage, from pachytene to zygotene, happens in the satellite-\emph{rich} subset
($\rho = 0.652$ against 0.639), a margin of 0.012, i.e. of the same order as the
Monte-Carlo noise of the 200-permutation null (0.005--0.02). We therefore report
the stage profile as descriptive
rather than as a tested difference: it is resolution-dependent (flat at
50\,kb, $\rho = 0.29$--0.34), it does not persist when satellite-rich windows
are removed, and the satellite-poor subset is 21 pairs drawn from 7 chromosome
ends, so its per-stage estimates are not independent. Minor-satellite content of
these windows could not be measured directly -- the local Dfam library has no
\emph{Mus} minor-satellite family, so MurSAT1, MMSAT4 and CENSAT\_MC were used
as a proxy, and they cover essentially none of the p-end flanks (mean fraction
$4\times10^{-5}$) -- and the curated major-satellite family GSAT\_MM covers
a mean of only 28.9\% of the p-end flank, so at least 59.6\% of it is covered by
the other seven Dfam families in the union.
%%% END INSERT (R1)

%%% BEGIN INSERT (R1-fallback) -- SHORTER, use if the stage ordering is
%%% confirmed indistinguishable after contact-level masking
Because the 1\,Mb mouse p-arm flanks are 26--99\% satellite (mean 88.5\%) and
satellite arrays are near-identical between chromosomes, these windows sit at
the top of the similarity axis and at the extreme of the multi-mapping problem
at once. The sequence--contact association survives when the satellite-rich
windows are set aside ($\rho = 0.486$--$0.539$ over the 21 satellite-poor arm
pairs, permutation $p \leq 0.025$), but the stage ordering does not
($\rho = 0.536$, 0.539, 0.486 and 0.486 across the four stages, no resolvable
maximum), so we report the mouse stage profile as descriptive only.
%%% END INSERT (R1-fallback)

%%% BEGIN INSERT (R2) -- summary paragraph, replaces submission/paper.tex:81--83
Mouse contact maps show the same preference:
the association strengthens as chromosomes synapse, although the stage ordering
is not resolved once the satellite-rich p-arm windows are set aside.
%%% END INSERT (R2)

%%% BEGIN INSERT (R3) -- Fig. 4c caption, replaces submission/paper.tex:545--547
%%% NOTE: the trailing "}" is the close of the \caption{ ... } group that is
%%% opened earlier in the original caption; it is part of the replaced excerpt,
%%% so this snippet is intentionally NOT brace-balanced on its own.
The arm-level Mantel test on
the same data, which respects that non-independence, gives $\rho = 0.41$, 0.58, 0.63 and
0.33 across the four stages (Methods); the mouse p-arm flanks are 88.5\%
satellite on average, and the stage ordering is descriptive only -- it does not
persist once the satellite-rich p-arm windows are set aside.}
%%% END INSERT (R3)
